\scanpage{12}\mspage{\handwritten{11}}

III \quad Point de vue ``motivique'' en théorie du cobordisme.

Soit $C$ la catégorie des variétés différe\unsure{b|n|m}tiables \inserted{\handwritten{(compactes?)}} (pas \unsure{nécessairement} \unsure{orinhtables|orientables}), \inserted{les morphismes étant les classes d'homotopie d'applications continues.} Si $B$ est une catégorie, on s'intéresse aux couples $(F_{\bullet},F^{\bullet})$ d'un foncteur covariant et d'un foncteur contravariant de $C$ dans $B$, satisfaisant les conditions que pour tout $X\handwritten{\in}\mathrm{Ob}\ C$, on a $F_{\bullet}(X)=F^{\bullet}(X)$, et que si on a un produit fibré ordinaire

% Four diagonal arrows drawn in pen between typed objects; labels are typed.
\diagram{\begin{tikzcd}[ampersand replacement=\&,row sep=small,column sep=small]
 \& X'\arrow[dl,"g'"']\arrow[dr,"f'"] \& \\
 X\arrow[dr,"f"'] \& \&Y'\arrow[dl,"g"]\\
 \&Y \& {,}
\end{tikzcd}}

avec $f$ et $g$ transversaux, alors on a $g^{\bullet}f_{\bullet}=\unsure{f!|f'_{\bullet}}\,g'^{\bullet}$ \inserted{(on pose $f_{\bullet}=F_{\bullet}(f),f^{\bullet}=F^{\bullet}(\unsure{f})$\unsure{)}}. La question est de trouver une $B$ et un couple $(F_{\bullet},F^{\bullet})$ qui soient ``universels'' dans un sens évident. Quillen prouve que la construction suivante donne une telle solution.

Pour toute variété $X\handwritten{\in}\mathrm{ob}\ C$, on définit un anneau commutatif $B(X)$ (l'anneau de bordisme de $X$) en prenant comme éléments les \unsure{classes} \deleted{de} d'éléments de $C$ au dessus de $X$, pour la relation de cobordisme: $f':X'\handwritten{\to}X$ et $f'':X''\handwritten{\unsure{-|\to}}X$ sont équivalents, s'il existe une variété à bord $Z$, et un $f:Z\handwritten{\to}X$, tel que $\mathrm{bord}\ Z=X'\handwritten{\amalg}X''$ et $f'$ et $f''$ sont induits par $f$. La somme de $B(X)$ est la somme disjointe (c'est une loi de \unsure{groupe} \deleted{l'inverse étant} où tout élément est d'ordre \deleted{1}\handwritten{2} \deleted{ou 2}), le produit est donné par le produit fibré, en utilisant le théorème de transversalité de Thom pour bouger un peu $f''$, de sorte à le rendre transversal à $f'$.

\handwritten{\uline{NB} L'anneau est de plus \uline{gradué} par la dimension relative....}

Si on a $f:X\handwritten{\to}Y$, on définit $f_{\bullet}:B(X)\handwritten{\to}B(Y)$ de façon évidente (\unsure{toute} variété sur $X$ est sur $Y$), et $f^{\bullet}:B(Y)\handwritten{\to}B(X)$ par produit fibré (utilisant encore le th. de transversalité). De cette façon, $B(X)$ devient un foncteur-anneau \inserted{\handwritten{(\unsure{gradué})}} contravariant en $X$, et un foncteur groupe covariant en

% archive sheet number: 11
